\begin{lemma}
\label{editorial-source-lemma-generating-finitely-generated-separable-field-extensions}
Let $K/k$ be a separably generated and finitely generated
field extension.
Set $r = \text{trdeg}_k(K)$. Then there exist elements
$x_1, \ldots, x_{r + 1}$ of $K$ such that
\begin{enumerate}
\item $x_1, \ldots, x_r$ is a transcendence basis of $K$ over $k$,
\item $K = k(x_1, \ldots, x_{r + 1})$, and
\item $x_{r + 1}$ is separable over $k(x_1, \ldots, x_r)$.
\end{enumerate}
\end{lemma}

\begin{proof}
Choose a separating transcendence basis of the
original extension. Since $K/k$ is finitely
generated, it has finite size $r$; enumerate
the chosen elements as $x_1,\ldots,x_r$.
Put $F=k(x_1,\ldots,x_r)$, retaining this
original subfield of $K$. If
$K=k(z_1,\ldots,z_m)$ is a finite generating
list, each $z_j$ is algebraic over $F$,
and $K=F(z_1,\ldots,z_m)$. The successive
finite algebraic degrees therefore prove
that $K/F$ is finite. It is separable by
the chosen separating basis. Fields,
Lemma \ref{fields-lemma-primitive-element}
supplies an original $x_{r+1}\in K$ with
$K=F(x_{r+1})$. Its minimal polynomial over
$F$ is separable because the extension is.
This proves all three assertions with the
same original elements. If $K=F$, take
$x_{r+1}=0$, whose minimal polynomial $T$
is separable. If $r=0$, the basis is empty,
$F=k$, and the same finite-extension
argument applies without any transcendental
generators.
\end{proof}